\begin{figure}[htbp]
\centering
\resizebox{\linewidth}{!}{%
\begin{tikzpicture}[font=\scriptsize]

\node[caja, text width=2.3cm, align=center, draw=ambar, line width=1pt] (sens) at (0,0)
  {\textbf{Sensores}\\ inalámbricos\\ a batería};
\node[caja, text width=2.3cm, align=center, draw=ambar, line width=1pt] (panel) at (4.6,0)
  {\textbf{Panel} o\\ comunicador};
\draw[flecha] (sens.east) -- node[nota, below=2pt, font=\scriptsize, align=center] {radio\\propietaria} (panel.west);

% los tres transportes, como cajas apiladas (ningún texto sobre líneas)
\node[caja, text width=4.9cm, align=left, draw=grislinea] (t1) at (9.0, 1.7)
  {\textbf{1. Línea telefónica fija}\\ tonos DTMF \ \textit{(histórico)}};
\node[caja, text width=4.9cm, align=left, draw=grislinea] (t2) at (9.0, 0)
  {\textbf{2. Red celular}\\ tonos DTMF sobre 2G/3G \ \textit{(en retirada)}};
\node[caja, text width=4.9cm, align=left, draw=verdeacento, line width=1.1pt, fill=verdeacento!6] (t3) at (9.0,-1.7)
  {\textbf{3. Internet}\\ SIA DC-09 con AES \ \textit{(el adoptado)}};

\node[caja, text width=2.4cm, align=center, draw=azuloscuro, line width=1.1pt, fill=azuloscuro!5]
  (arc) at (14.4,0) {\textbf{Central}\\ \textbf{receptora}\\ (ARC / CMS)};

\foreach \t in {t1,t2,t3}{
  \draw[flecha] (panel.east) -- (\t.west);
  \draw[flecha] (\t.east) -- (arc.west);
}

\node[notaclave, text width=14.2cm, align=left, font=\footnotesize] at (7.2,-4.0)
  {\textbf{En los tres casos el mensaje de aplicación es el mismo: Contact ID.} La «comunicación celular»
   es un cambio de transporte, no de protocolo. El mensaje tiene 16 dígitos porque debía caber en tonos
   telefónicos sobre una línea analógica, \textbf{y esa restricción de hace cuatro décadas es la que hoy
   lo hace apropiado para LoRa}, un canal igualmente estrecho.};
\end{tikzpicture}}
\caption{Cadena de comunicación de una alarma monitoreada: el formato del mensaje se conserva a través
de tres generaciones de transporte.}
\label{fig:alarmas}
\end{figure}
